
\documentclass{article}
\usepackage{colortbl}
\usepackage{makecell}
\usepackage{multirow}
\usepackage{supertabular}

\begin{document}

\newcounter{utterance}

\centering \large Interaction Transcript for game `cladder', experiment `full\_v1.5\_default', episode 6328 with score{-}Qwen3.5{-}9B{-}sp.
\vspace{24pt}

{ \footnotesize  \setcounter{utterance}{1}
\setlength{\tabcolsep}{0pt}
\begin{supertabular}{c@{$\;$}|p{.15\linewidth}@{}p{.15\linewidth}p{.15\linewidth}p{.15\linewidth}p{.15\linewidth}p{.15\linewidth}}
    \# & \multicolumn{2}{c}{Player} && \multicolumn{2}{c}{Game Master} \\
    \hline

    \theutterance \stepcounter{utterance}  
    & & & \multicolumn{4}{p{0.6\linewidth}}{
        \cellcolor[rgb]{0.9,0.9,0.9}{
            \makecell[{{p{\linewidth}}}]{
                \texttt{\tiny{[P1$\langle$GM]}}
                \texttt{You are an expert in causal inference. The following question is not a typical commonsense query, but rather a meticulously designed question created by a professor specializing in causal inference, intended to assess the students' mastery of the course content.} \\
\texttt{Imagine a self{-}contained, hypothetical world with only the following conditions, and without any unmentioned factors or causal relationships: Swoy has a direct effect on rixq and xevu. Rixq has a direct effect on zuph. Zuph has a direct effect on xevu. Swoy is unobserved. For those who are not rixq and are not zuph, the probability of xevu is 46\%. For those who are not rixq and are zuph, the probability of xevu is 41\%. For those who are rixq and are not zuph, the probability of xevu is 66\%. For those who are rixq and are zuph, the probability of xevu is 60\%. For those who are not rixq, the probability of zuph is 61\%. For those who are rixq, the probability of zuph is 34\%. For those who are rixq, would it be more likely to see xevu if the individual was not rixq?} \\
\texttt{Start your answer with 'Yes' or 'No', followed by additional reasoning or evidence to support your explanation.} \\
            }
        }
    }
     \\ \\

    \theutterance \stepcounter{utterance}  
    & \multicolumn{4}{p{0.6\linewidth}}{
        \cellcolor[rgb]{0.561,0.737,0.561}{
            \makecell[{{p{\linewidth}}}]{
                \texttt{\tiny{[P1$\rangle$GM]}}
                \texttt{No} \\
\\ 
\texttt{**Reasoning:**} \\
\\ 
\texttt{To answer this question, we must determine the causal effect of the variable **rixq** on **xevu**. The question asks: "For those who are rixq, would it be more likely to see xevu if the individual was not rixq?"} \\
\\ 
\texttt{This is asking for the counterfactual comparison: \$P(xevu | do(rixq=1))\$ vs. \$P(xevu | do(rixq=0))\$. Specifically, we are comparing the probability of observing \$xevu\$ given that the individual *is* \$rixq\$ (the factual state) against the probability of observing \$xevu\$ if that same individual *were not* \$rixq\$ (the counterfactual state).} \\
\\ 
\texttt{**Step 1: Analyze the Causal Structure**} \\
\texttt{The problem states:} \\
\texttt{1.  Swoy \$\textbackslash{}rightarrow\$ rixq} \\
\texttt{2.  Swoy \$\textbackslash{}rightarrow\$ xevu} \\
\texttt{3.  rixq \$\textbackslash{}rightarrow\$ zuph} \\
\texttt{4.  zuph \$\textbackslash{}rightarrow\$ xevu} \\
\\ 
\texttt{We are given the following probabilities:} \\
\texttt{*   \$P(xevu | \textbackslash{}neg rixq, \textbackslash{}neg zuph) = 0.46\$} \\
\texttt{*   \$P(xevu | \textbackslash{}neg rixq, zuph) = 0.41\$} \\
\texttt{*   \$P} \\
            }
        }
    }
    & & \\ \\

    \theutterance \stepcounter{utterance}  
    & & & \multicolumn{2}{p{0.3\linewidth}}{
        \cellcolor[rgb]{0.9,0.9,0.9}{
            \makecell[{{p{\linewidth}}}]{
                \texttt{\tiny{[GM$|$GM]}}
                \texttt{No} \\
            }
        }
    }
    & & \\ \\

    \theutterance \stepcounter{utterance}  
    & & & \multicolumn{2}{p{0.3\linewidth}}{
        \cellcolor[rgb]{0.9,0.9,0.9}{
            \makecell[{{p{\linewidth}}}]{
                \texttt{\tiny{[GM$|$GM]}}
                \texttt{no} \\
            }
        }
    }
    & & \\ \\

    \theutterance \stepcounter{utterance}  
    & & & \multicolumn{2}{p{0.3\linewidth}}{
        \cellcolor[rgb]{0.9,0.9,0.9}{
            \makecell[{{p{\linewidth}}}]{
                \texttt{\tiny{[GM$|$GM]}}
                \texttt{game\_result = WIN} \\
            }
        }
    }
    & & \\ \\

\end{supertabular}
}

\end{document}
